\section{Experimental Design}\label{sec:experiments}

The original confirmatory protocol was fixed before observing its outcomes.
A compact registry records its method version, cells, failures, and analysis
endpoints.  Supplementary comparisons added after manuscript review reuse
those target seeds and are identified separately below.  Every implementable
method sees target truth only after its deployment decision, when truth is
used to compute evaluation metrics.

\subsection{Randomized ordered-profile stress test}

The primary study uses eight fixed, equally weighted regimes.  Within each
regime, independent deterministic seed streams generate held-out problem
instances.  Table \ref{tab:stress-regimes} separates grid resolution from the
complexity of the underlying policy shape.


For each regime, we draw 20 independently seeded latent target tasks and cross
the same seeds with $d\in\{200,1000,10000\}$.  This gives 160 independent
latent task seeds and 480 correlated task--resolution cells per design.  Each
cell in the original frozen matrix is paired across six designs: source-scored,
unstandardized DCT
maximin, random low frequency, natural three-block, raw Sobol, and an
explicitly labeled finite-library upper bound.  The primary implementable
matrix has
$8\times3\times20\times5=2400$ cells; including the oracle gives 2880.  Every
design uses $n_0=N=10$, so the primary comparison is a pure initial-design
intervention.  The source-scored arm pays 384 source calls; source-free arms
pay none.  All use the same independent shortlist verifier.

\paragraph{Supplementary matched comparisons.}\label{sec:supplemental-controls}
We added standardized DCT maximin and the source-greedy portfolio on the same
160 latent tasks at all three resolutions.  These controls use $n_0=N=10$,
the original target noise streams, and the primary $(80,80,80)$ verifier.
Standardized DCT matches the source-scored atlas's structural scaling; the
portfolio also matches its 384-call source archive.  The standardized control
at $d=1000$ was first evaluated during manuscript review and is reused
unchanged; the remaining runs followed the recorded supplementary rules.
All rules were chosen after the original target results had been examined.
They isolate the comparison confound and
selection-rule choice on reused tasks; they are not a new independent
confirmation.  The original matrix and its results remain frozen.

For each synthetic target regime, the two supplied source tasks are generated
from that same regime with the role set to source.  This is a regime-matched
transfer experiment, not a source-retrieval experiment.  Algorithms receive
neither target outcomes nor hidden target geometry, but the experimenter has
already supplied a related archive.  The Energy region holdout is the explicit
archive-mismatch control.

The primary chance level is $\alpha=.05$, and task safe mass is calibrated to
.08 on a separate fixed library that is not used for design.  Regret is
measured against the best feasible policy in that evaluation library.  The
$d=10000$ cells refine the grid while holding the underlying problem fixed
unless the regime explicitly changes rank.

\paragraph{First-center attribution and new-family confirmation.}
We subsequently fixed a factorial comparison of the structural medoid and
source-best first center, crossed with structural-only and augmented metrics.
The medoid was held fixed across the two metrics.  A fifth design used the
source-best center with Voronoi-weighted raw-profile $L^2$ distance on the same
128 reference nodes.  These diagnostic comparisons reused the original tasks;
identical ordered designs reused their original outcomes with the proper
source cost.  A separate local prospective protocol then fixed three new
family seeds, three corresponding public libraries and source archives, and
ten new tasks per mechanism and family at $d=1000$.  Its 240 latent tasks
were compared across six designs, including the same-source portfolio.
The primary contrast was source scoring minus standardized generic maximin
in certified true-feasible deployment.  Intervals resampled tasks within the
24 fixed family--mechanism strata, and each family was reported separately.
Library sizes 32 and 128 were additional fixed sensitivities on the first new
family; they change support composition as well as cardinality.  The archive,
selection rule and verifier were not tuned after these new outcomes.

\subsection{Sensitivity and information controls}

An 8640-cell one-factor-at-a-time matrix varies active rank, $\alpha$, safe
mass, $n_0$, number of source tasks, profiles per task, source replications,
retained frequency, frequency penalty, source-rank metric weights, and the
first-center safety weight.  These experiments describe sensitivity; they are
not a hyperparameter search performed after seeing the primary results.  A
separate 1280-cell matrix crosses declared versus
schema-blind profile mapping with equal versus descriptor-conditioned source
weights.  The latter weighting is excluded from the primary method regardless
of its result.

\subsection{Cost and optimizer controls}

The equal-preverification-cost matrix compares 384 source plus 10 target calls
with 394 target-only calls for each source-free design.  The same independent
verifier is applied afterward and all-in realized calls include its cost.

Target-only functional SCBO searches eight nonconstant cosine coefficients
and one level parameter (nine optimization variables) with the
published constrained trust-region posterior-sampling backend
\citep{BalandatEtAl2020,ErikssonPoloczek2021}.  We test coefficient counts
4, 8, and 16 at $N=13$, and the primary 8-coefficient version at $N=394$.
One $N=394$ task in the irregular-grid regime returned duplicate canonical
posterior samples; it remains a failed, unsuccessful cell.

Eight native transfer pipelines (Safe F-PACOH, RGPE, stacked transfer GP,
multi-task GP, FSBO, HyperBO, MetaBO, and MALIBO) receive the same 384-call
source archive, $n_0=10$, and $N=13$, but retain their own
prior, representation, initialization, and acquisition logic.  This is an
end-to-end transfer comparison on three fixed families with 20 algorithmic
seeds per family.  All eight pipelines share the historical v69 terminal
protocol: candidate budgets $(80,128,128)$ and an 8-call incumbent objective
comparison.  This protocol differs from the primary $(80,80,80)$ verifier;
native results are interpreted within their own matrix.

\subsection{Electricity-reserve external comparison}

The external study uses the Open Power System Data time-series release
\citep{OPSD2020}.  A 1000-point policy maps observable forecast stress to
reserve responses over a 168-hour service window.  The corrected Energy V5
model gives every policy the same initial SOC, 50\% of capacity, charges all
subsequent grid purchases, and records terminal energy.  Reserve adjustment
and forecast-error balancing share one hourly converter budget: the sum of
external charging and discharging energy cannot exceed the power capacity
times one hour.  Reserve adjustment is executed first and balancing uses the
remaining budget.  Its objective is
finite-period service cost with zero terminal salvage value, rather than a
steady-state economic criterion.

All eight designs were rerun on 18 markets in five geographic regions with
five algorithmic seeds per market, giving 720 cells.  Each uses $n_0=10$,
$N=13$, and the same $(80,80,80)$ verifier.  One source archive is frozen for
each excluded geographic region before target evaluation, with four external
source markets, 64 profiles and three replications (768 calls).  Its source
outcomes and selected designs are shared across the region's target markets
and source-based comparisons.  The data splits remain search in 2017, audit
in 2018 and verification in 2019.  V2--V4 results are preserved as
historical evidence; V5 supplies the current external comparison.

Fresh verification windows are sampled from the held-out empirical calendar.
After each decision, four chronological blocks and physically nonoverlapping
windows provide descriptive stability checks.  Actual reuse accounting charges
one archive per excluded region and sums target costs across its observed
markets; five seeds estimate repeatability within each market and do not count
as five deployments.  This corrected study reuses previously examined data,
so it is distinct from the prospective synthetic confirmation.

\subsection{Inference and failure accounting}

For randomized profiles, a latent task seed within a fixed regime and
resolution is the generalization unit.  Supplementary paired bootstrap
intervals resample task seeds within each fixed regime and retain the equal
regime weights.  Task-law concentration is applied separately at each resolution
with regimes as fixed strata.  Rates pooled over all three resolutions
describe 480 task--resolution cells but are not treated as 480 independent
tasks.  Simulation seeds within a fixed external market measure algorithmic
repeatability, while geographic
regions are the conservative Energy comparison units.  We do not pool seeds
as if they were independent domains.  Exact sign tests are Holm-adjusted within
their registered families, but effect sizes and intervals receive priority
over very small $p$-values.  Every timeout, duplicate-sampling failure, and
false certificate remains in its denominator.

Primary endpoints are feasible coverage, certified true-feasible deployment,
false certification, and penalized loss.  We additionally report source,
search, verification, unamortized all-in, and $M=20$ amortized calls; conditional
regret is secondary.
